% Round 17. A non-circular global existence construction.
% Depends on independently reviewed ON17 ordinary-node activation and PH16.
\section{Zariski-dense hypersurface families with full local G and BB height}
\label{r17gp:section}

\paragraph{Status and scope.}
The construction below produces specified families over $\mathbf Q$,
and then over every number field by base change. It does not assert
that every variety belongs to these families. The local input is
Theorem~\ref{r17on:actual}: good reduction, or a geometrically integral
special fiber with isolated ordinary quadratic singularities in odd
residue characteristic, gives ordinary actual strong G in all degrees.
That input allows arbitrary finite thickness at the nodes. Its exact
native normal form is Proposition~\ref{r17on:normal-form}, and its
proof uses QF17 and the frozen actual detector. The global height input
is the frozen PH16 theorem \texttt{r16ph:general-height}.

The classical geometric theory of Lefschetz pencils is a published
input in the literature; see Jannsen--Saito,
\emph{Bertini theorems and Lefschetz pencils over discrete valuation
rings}, \url{https://arxiv.org/pdf/0911.1470}, and
Esnault--Kerz, \emph{Semistable Lefschetz pencils},
\url{https://page.mi.fu-berlin.de/esnault/preprints/helene/151_crelle.pdf}.
We give the specific characteristic-zero fixed-member argument used
here, so that no unstated freedom to change that member enters the
arithmetic construction. No novelty claim is made.

Fix $n\ge3$, put $N=n+1$, and let $D\ge3$. Write
\[
 F_0=X_0^D+\cdots+X_N^D,\qquad
 J=\sum_{i=0}^{N-1}X_i^{D-1}X_{i+1}+X_N^D,
 \qquad b_D=\prod_{p\mid D}p.
 \tag{GP17.1}\label{r17gp:polynomials}
\]
The symbol $b_D$ is used only for this fixed integer; it is not a
node thickness.

\begin{lemma}[A smooth member at every prime dividing the degree]
\label{r17gp:chain}
For every prime $p\mid D$, the hypersurface $J=0$ in
$\mathbf P^N_{\mathbf F_p}$ is smooth. Therefore, for any integral
form $F_1\equiv J\pmod{b_D}$, at every prime $p$ at least one
of $F_0=0$ and $F_1=0$ has smooth reduction: use $F_0$ if
$p\nmid D$, and $F_1$ if $p\mid D$.
\end{lemma}

\begin{proof}
Work over an algebraic closure of $\mathbf F_p$. If all partial
derivatives of $J$ vanish, its last derivative gives
\[
                        X_{N-1}^{D-1}=0.
\]
Thus $X_{N-1}=0$. For $1\le i\le N-1$,
\[
 \frac{\partial J}{\partial X_i}
    =X_{i-1}^{D-1}+(D-1)X_i^{D-2}X_{i+1}.
\]
Descending induction gives $X_{N-2}=\cdots=X_0=0$.
The use of $X_i^{D-2}=0$ is valid because $D\ge3$.
A projective point with these coordinates has $X_N\ne0$,
whereas $J=X_N^D\ne0$. Hence no point of $J=0$ is singular.
If $p\nmid D$, the Fermat gradients $D X_i^{D-1}$ cannot
all vanish at a projective point. Reduction of the stated
congruence proves the assertion for $F_1$.
\end{proof}

\begin{lemma}[Lefschetz pencils through a prescribed smooth hypersurface]
\label{r17gp:fixed-member}
Let $k$ be an infinite field of characteristic zero, let
$N\ge2$, $D\ge3$, and let $f_0$ define a smooth degree-$D$
hypersurface in $\mathbf P^N_k$. There is a nonempty Zariski
open set of degree-$D$ forms $f_1$ such that:
\begin{enumerate}
\item $f_0=f_1=0$ is a smooth codimension-two complete intersection;
\item $[u:v]\mapsto\{u f_0+v f_1=0\}$ is a Lefschetz pencil:
      every singular geometric member has exactly one ordinary
      quadratic singularity, and its line meets the discriminant
      transversely there;
\item the member $f_1=0$ is smooth.
\end{enumerate}
\end{lemma}

\begin{proof}
It suffices to construct the open set geometrically and observe
that each condition descends to $k$. Let $\mathbf P^M$ be the
space of degree-$D$ forms. The incidence
\[
 I=\{(f,x): f(x)=0,\ d f(x)=0\}
                      \subset\mathbf P^M\times\mathbf P^N
\]
uses the first jet of a local trivialization of $\mathcal O(D)$.
Evaluation of first jets is surjective. Thus $I$ is a smooth
irreducible projective bundle over $\mathbf P^N$ of dimension
$M-1$.

The degeneracy of the quadratic term at the marked singular
point is one proper closed condition on $I$: in coordinates
$x=[1:0:\cdots:0]$, forms containing
$X_0^{D-2}(X_1^2+\cdots+X_N^2)$ have nondegenerate quadratic
term. The incidence of a form with two distinct marked singular
points has dimension at most $M-2$. Indeed first-jet evaluations
at two distinct points are jointly surjective for $D\ge3$.
After changing coordinates to the first two coordinate points,
the $2(N+1)$ monomials specifying those first jets are distinct;
this verifies independence directly. Its fibers over the ordered
pair of points have the predicted projective dimension, and the
pair space has dimension $2N$.

Consequently a dense open of $I$ consists of forms with precisely
one singular point and a nondegenerate quadratic term. Its image
has dimension $M-1$, so the discriminant $\Delta\subset\mathbf P^M$
is an irreducible hypersurface. Outside a closed subset of
$\mathbf P^M$ of dimension at most $M-2$, its points are exactly
these one-node forms. The discriminant is smooth there: the
invertible Hessian lets the critical-point equations be solved
uniquely etale locally as the coefficients vary, and the remaining
critical-value equation has nonzero differential, given by
evaluation of the coefficient variation at the marked point.
Denote this open by $\Delta^\circ$.

Project from the fixed point $[f_0]\notin\Delta$ to the space
$\mathbf P^{M-1}$ of lines through it. Restricted to $\Delta$
this is finite: a fiber is a line intersection, and a line through
$[f_0]$ cannot be contained in $\Delta$. It is generically etale,
because the characteristic is zero. Removing its branch locus
and the image of $\Delta\setminus\Delta^\circ$ gives a nonempty
open set of directions whose lines meet $\Delta$ only at smooth
one-node points and transversely.

There is also a nonempty open set of directions for which the
base intersection with $f_0=0$ is smooth. This is ordinary Bertini
for the globally generated very ample restriction of
$\mathcal O(D)$ to the fixed smooth hypersurface. Its section
space is supplied by the restrictions of ambient forms. Intersect
the two opens in the irreducible direction space. Lifting to the
space of choices of $f_1$ and excluding $\Delta$ itself gives the
asserted nonempty open set. Smoothness of the base intersection
also ensures that every pencil member is smooth along the base:
the two ambient differentials are independent there.
\end{proof}

\begin{lemma}[An integral pencil with a good member at every prime]
\label{r17gp:integral-pencil}
There exists an integral degree-$D$ form
\[
                              F_1\equiv J\pmod{b_D}
 \tag{GP17.2}\label{r17gp:congruence}
\]
such that the pencil through $F_0,F_1$ is as in
Lemma~\ref{r17gp:fixed-member}. The set of such integral forms
is Zariski dense in the affine space of degree-$D$ forms over
$\mathbf Q$.

For each such choice there is a finite set $S$ of rational
primes, containing every prime dividing $2D$, such that the
pencil over $\mathbf Z[1/S]$ has smooth total space and smooth
base locus, and every geometric member is smooth or has exactly
one ordinary quadratic singularity. Its critical scheme is
finite etale over $\mathbf Z[1/S]$ and maps as a closed immersion
into $\mathbf P^1_{\mathbf Z[1/S]}$.
\end{lemma}

\begin{proof}
The affine lattice of coefficient tuples
$J+b_D\mathbf Z^{M+1}$ is Zariski dense over $\mathbf Q$.
For completeness, a polynomial vanishing on a product of
infinite arithmetic progressions is zero by induction on the
number of variables, using the one-variable root bound at
each step. The same observation shows that its intersection
with any nonempty open set is Zariski dense in that open.
Apply Lemma~\ref{r17gp:fixed-member} to the Fermat form.

Over $\mathbf Q$ the incidence hypersurface
\[
 \mathcal P=\{uF_0+vF_1=0\}
                  \subset\mathbf P^N\times\mathbf P^1
\]
is smooth. Away from the base, a derivative in the parameter
direction is nonzero; at the base, independence of $dF_0,dF_1$
gives smoothness. The family is flat over $\mathbf P^1$ because
all its fibers are nonzero degree-$D$ hypersurfaces. Its critical
points are nondegenerate, form a finite reduced scheme, and have
distinct geometric critical values. In characteristic zero this
scheme is finite etale, and its map to $\mathbf P^1$ is a closed
immersion.

All these data and equations are of finite presentation over
$\mathbf Z$. After inverting finitely many primes they retain
these properties. More explicitly one removes primes of failure
of the smooth base intersection or smooth total space, removes
vertical components of the relative critical locus, and then
spreads its finite etale structure, its closed-immersion map to
the parameter line, and the invertible Hessian. These are open
conditions or vanishing conditions on finitely presented
modules which already hold over $\mathbf Q$. We also remove
primes at which the two coefficient vectors become dependent,
so that no parameter gives the zero polynomial. Adding the primes
dividing $2D$ gives the asserted finite set $S$. The closed
relative critical locus has no other points, so the Hessian
condition proves the stated description of every geometric
fiber. This spreading step makes no claim that $S$ is empty.
\end{proof}

\begin{lemma}[Geometric integrality of an isolated-node hypersurface]
\label{r17gp:integral-fiber}
A degree-$D$ hypersurface in $\mathbf P^{n+1}$ over an algebraically
closed field, with $n\ge2$ and with at most isolated singular
points, is integral.
\end{lemma}

\begin{proof}
A repeated irreducible factor makes an entire component singular,
which is impossible. If there are two distinct components, their
intersection in projective space has dimension at least $n-1$,
is nonempty, and is contained in the singular locus of their
product equation. Since $n-1\ge1$, this too is impossible.
Thus the defining equation is irreducible and reduced. Applying
this over the algebraic closure proves the corresponding
geometric assertion. Smooth members are covered by the same
argument, or by connectedness of a positive-degree hypersurface
of positive dimension and regularity.
\end{proof}

\begin{theorem}[Infinitely many full-height members over the original field]
\label{r17gp:global-family}
Fix an integral $F_1$ and finite set $S$ as in
Lemma~\ref{r17gp:integral-pencil}. Put $M_D=D$ (any fixed
positive power of $D$ would also work). There are infinitely
many integers $a$ for which
\[
 X_a=\{M_DF_0+aF_1=0\}\subset\mathbf P^{n+1}_{\mathbf Q}
 \tag{GP17.3}\label{r17gp:member}
\]
is smooth geometrically integral and has the following
properties.
\begin{enumerate}
\item At every finite place of $\mathbf Q$, every specified
      geometrically homologically trivial rational cycle on
      $X_a$ satisfies ordinary actual strong G in every
      codimension, with its exact specified support and all
      original actual compact tests.
\item The same assertion holds after base change to every
      number field $F$, at every finite place of $F$, for
      all cycles defined over that completion, including
      cycles not descending to $\mathbf Q$.
\item For every number field $F$ and every $1\le r\le n$,
      PH16 gives the canonical real-valued Beilinson--Bloch pairing
      on the full geometric homological Chow groups
      \[
       \operatorname{CH}^r((X_a)_F)_{\mathbf Q,\mathrm{geom\,hom}}
       \ \times\ 
       \operatorname{CH}^{n+1-r}((X_a)_F)_{\mathbf Q,\mathrm{geom\,hom}}
                           \longrightarrow\mathbf R.
       \tag{GP17.4}\label{r17gp:height}
      \]
\end{enumerate}
No squarefreeness condition on a discriminant value is imposed.
The statement gives infinitely many parameter members; pairwise
nonisomorphism is not part of the assertion.
\end{theorem}

\begin{proof}
Choose $a$ in the arithmetic progression specified by
\[
        a\equiv1\pmod{b_D},\qquad
        a\equiv0\pmod p\quad(p\in S,\ p\nmid D).
 \tag{GP17.5}\label{r17gp:parameter-crt}
\]
The moduli are relatively prime after collecting distinct
primes, so the Chinese remainder theorem gives an infinite
progression. Exclude its finitely many parameters whose
characteristic-zero members lie on the discriminant of the
fixed pencil. All remaining members are smooth. Their geometric
integrality follows from Lemma~\ref{r17gp:integral-fiber}.

If $p\mid D$, the integral equation in
\eqref{r17gp:member} reduces to a nonzero scalar multiple
of $F_1$, and is smooth by Lemma~\ref{r17gp:chain}.
If $p\in S$ and $p\nmid D$, it reduces to a nonzero scalar
multiple of $F_0$, again smooth. These equations give proper
smooth models over $\mathbf Z_p$: the equations are primitive,
hence flat, and the Jacobian criterion on the special fiber
gives smoothness there; the generic fibers are already smooth.
In particular every prime of residue characteristic two has
been placed in this good-reduction part.

For $p\notin S$, the coefficient $M_D$ is a unit. Thus the
parameter is the integral point $[M_D:a]$ of the spread pencil.
Its special fiber is smooth or has one ordinary quadratic
singularity, and is geometrically integral by
Lemma~\ref{r17gp:integral-fiber}. The equation is a nonzero
hypersurface equation in each fiber and gives a flat projective
model. At a singular point Theorem~\ref{r17on:normal-form}
puts its henselian, equivalently pointed-etale, equation into
native form
\[
                             Q(x)=u\pi^b,
                       \qquad b\ge1,
 \tag{GP17.6}\label{r17gp:thickness}
\]
with $Q$ nondegenerate and $u$ a unit. The thickness $b$ is
finite because the generic fiber is smooth. There is no upper
bound or squarefreeness restriction on $b$; it is precisely
this arbitrary-thickness case which ON17 and QF17 handle.
Theorem~\ref{r17on:actual} proves full local actual G.
At the good places the frozen good-reduction theorem does so.
This proves the first assertion.

Now let $F$ be any number field and $w\mid p$ any finite
place. Good reduction stays good after the base change
$\mathbf Z_p\to\mathcal O_{F_w}$. At the remaining places,
the special fiber remains geometrically integral and its
ordinary nodes remain ordinary after residue extension.
In the displayed local equation the thickness is multiplied
by the ramification index $e(F_w/\mathbf Q_p)$, with the
remaining factor absorbed into the unit. Equivalently one
applies the native normal-form theorem anew over
$\mathcal O_{F_w}$. The residue characteristic is still odd,
so the same full local theorem applies to every geometrically
homologically trivial rational cycle on $(X_a)_{F_w}$.
This is a new application of that local theorem over $F_w$,
not merely trace descent or pullback for cycles originally
defined over $\mathbf Q_p$. It proves the second assertion.

All finite places and both complementary codimensions now
satisfy the exact strong-G hypothesis of PH16. Apply its
full-Chow general-G height bridge with the normalizations
over the original number field $F$. This gives
\eqref{r17gp:height}. No assertion about other varieties is
used in this last step.
\end{proof}

\begin{corollary}[Zariski density and ample canonical bundle]
\label{r17gp:dense}
For each $n\ge3$ and $D\ge3$, the smooth degree-$D$
hypersurfaces over $\mathbf Q$ obtained above, as $F_1$
and $a$ vary, are Zariski dense in the parameter space of
smooth degree-$D$ hypersurfaces in $\mathbf P^{n+1}$.
Every one of them has the full height pairing
\eqref{r17gp:height} over every number field after base change.
If $D\ge n+3$, each such hypersurface has ample canonical
bundle.
\end{corollary}

\begin{proof}
For each allowed $F_1$, infinitely many parameters on its
fixed pencil occur in Theorem~\ref{r17gp:global-family}.
They are Zariski dense on that projective line. Therefore
any closed subset of the full coefficient parameter space
containing all constructed members contains that line and,
in particular, the point $[F_1]$. The allowed coefficient
tuples $F_1\equiv J\pmod{b_D}$ lying in the fixed-member
Lefschetz open are Zariski dense by
Lemma~\ref{r17gp:integral-pencil}. Hence the closed subset
must be the entire parameter space. Intersecting with its
smooth open gives the assertion.

Adjunction gives
$K_{X_a}=\mathcal O_{X_a}(D-n-2)$, which is ample when
$D\ge n+3$.
\end{proof}

\paragraph{What this does and does not settle.}
The result supplies unrestricted homological Chow-group heights
for Zariski-dense families, including dimension four and arbitrarily
high dimension, and allows arbitrary local node thickness.
It does not prove that every smooth hypersurface, or every smooth
projective variety, can be moved into the family while preserving
its original Chow classes. Zariski density alone gives no such
specialization or deformation theorem. General strong G and the
full unconditional height construction for an arbitrary variety
therefore remain open outside the independently established
classes. The local and global degree calibrations, exact support,
and full actual testing requirements are inherited from ON17,
QF17, and PH16 without weakening them.
